% Ngân hàng bài dummy (tự sinh để test): Hàm số y = ax^2 (a khác 0) (lớp 9).
% Quy ước:
%   - Mỗi bài: \begin{problem}[id] ... \end{problem}; lời giải lồng bên trong.
%   - id: duy nhất, không dấu, không đổi, không dùng lại. Dạng <lớp><phân môn>-<chủ đề>-<số>.
%   - Mỗi bước lời giải là 1 đoạn, các bước cách nhau 1 dòng trống; mỗi bước có lý do.
%   - Dòng cuối lời giải: Đáp số: $\boxed{...}$
%   - Cách giải khác: thêm \begin{solution}[cách 2].
%   - Metadata bằng comment ngay sau \begin{problem}: % nguon:, % lop:, % kc_du_kien:, % do_kho_du_kien:
%   - Đồ thị được mô tả bằng lời, không vẽ hình.
\documentclass[12pt,a4paper]{article}
\usepackage[utf8]{vietnam}
\usepackage{amsmath,amssymb}
\newenvironment{problem}[1][]{\par\medskip\noindent\textbf{Bài #1.}\ }{\par}
\newenvironment{solution}[1][]{\par\noindent\textit{Lời giải #1.}\par}{\par}

\begin{document}

% ---------------------------------------------------------------- tính giá trị hàm số
\begin{problem}[9DS-HSB2-001]
% nguon: dummy (tự sinh để test)
% lop: 9
% kc_du_kien: 9DS-HSB2-01
% do_kho_du_kien: 1
Cho hàm số $y = f(x) = 2x^2$.

a) Tính $f(-3)$, $f\left(\dfrac{1}{2}\right)$ và $f\left(\sqrt{2}\right)$.

b) Tìm các giá trị của $x$ để $y = 8$.

\begin{solution}
a) Thay $x = -3$ vào công thức của hàm số: $f(-3) = 2 \cdot (-3)^2 = 2 \cdot 9 = 18$.

a) Thay $x = \dfrac{1}{2}$: $f\left(\dfrac{1}{2}\right) = 2 \cdot \left(\dfrac{1}{2}\right)^2 = 2 \cdot \dfrac{1}{4} = \dfrac{1}{2}$.

a) Thay $x = \sqrt{2}$: $f\left(\sqrt{2}\right) = 2 \cdot \left(\sqrt{2}\right)^2 = 2 \cdot 2 = 4$.

b) $y = 8 \Leftrightarrow 2x^2 = 8 \Leftrightarrow x^2 = 4$ (chia hai vế cho $2$).

b) Hai số có bình phương bằng $4$ là $2$ và $-2$, nên $x = 2$ hoặc $x = -2$.

Đáp số: a) $\boxed{f(-3) = 18,\ f\left(\dfrac{1}{2}\right) = \dfrac{1}{2},\ f\left(\sqrt{2}\right) = 4}$; b) $\boxed{x = 2;\ x = -2}$
\end{solution}
\end{problem}

% ---------------------------------------------------------------- tìm a, điểm thuộc đồ thị
\begin{problem}[9DS-HSB2-002]
% nguon: dummy (tự sinh để test)
% lop: 9
% kc_du_kien: 9DS-HSB2-01
% do_kho_du_kien: 2
Cho hàm số $y = ax^2$ $(a \neq 0)$.

a) Tìm $a$, biết đồ thị của hàm số đi qua điểm $A(-2; 2)$.

b) Với giá trị $a$ vừa tìm được, các điểm $B(4; 8)$, $C\left(-1; -\dfrac{1}{2}\right)$, $D\left(\sqrt{6}; 3\right)$ có thuộc đồ thị của hàm số hay không? Vì sao?

\begin{solution}
a) Đồ thị đi qua $A(-2; 2)$ nên toạ độ của $A$ thoả mãn công thức hàm số: $2 = a \cdot (-2)^2 = 4a$.

a) Suy ra $a = \dfrac{1}{2}$ (thoả mãn $a \neq 0$); hàm số là $y = \dfrac{1}{2}x^2$.

b) Điểm $B(4; 8)$: thay $x = 4$ được $y = \dfrac{1}{2} \cdot 16 = 8$, bằng tung độ của $B$, nên $B$ thuộc đồ thị.

b) Điểm $C\left(-1; -\dfrac{1}{2}\right)$: thay $x = -1$ được $y = \dfrac{1}{2} \cdot 1 = \dfrac{1}{2} \neq -\dfrac{1}{2}$, nên $C$ không thuộc đồ thị (cũng thấy được vì $a > 0$ nên mọi điểm của đồ thị có tung độ không âm).

b) Điểm $D\left(\sqrt{6}; 3\right)$: thay $x = \sqrt{6}$ được $y = \dfrac{1}{2} \cdot \left(\sqrt{6}\right)^2 = \dfrac{1}{2} \cdot 6 = 3$, bằng tung độ của $D$, nên $D$ thuộc đồ thị.

Đáp số: a) $\boxed{a = \dfrac{1}{2}}$; b) $\boxed{B, D \text{ thuộc đồ thị; } C \text{ không thuộc đồ thị}}$
\end{solution}
\end{problem}

% ---------------------------------------------------------------- tính chất, so sánh, điểm có tung độ cho trước
\begin{problem}[9DS-HSB2-003]
% nguon: dummy (tự sinh để test)
% lop: 9
% kc_du_kien: 9DS-HSB2-02, 9DS-HSB2-01
% do_kho_du_kien: 2
Cho hàm số $y = f(x) = -\dfrac{1}{3}x^2$.

a) Hàm số đồng biến, nghịch biến khi nào? Mô tả đồ thị của hàm số (đỉnh, trục đối xứng, vị trí so với trục hoành).

b) Không tính giá trị, hãy so sánh $f\left(2\sqrt{2}\right)$ và $f(3)$.

c) Tìm các điểm thuộc đồ thị của hàm số có tung độ bằng $-3$.

\begin{solution}
a) Hàm số có dạng $y = ax^2$ với $a = -\dfrac{1}{3} < 0$.

a) Vì $a < 0$ nên hàm số đồng biến khi $x < 0$ và nghịch biến khi $x > 0$.

a) Đồ thị là một parabol có đỉnh $O(0; 0)$, nhận trục $Oy$ làm trục đối xứng; vì $a < 0$ nên đồ thị nằm phía dưới trục hoành và $O$ là điểm cao nhất của đồ thị.

b) Ta có $\left(2\sqrt{2}\right)^2 = 8 < 9 = 3^2$ và $2\sqrt{2}$, $3$ đều dương, nên $0 < 2\sqrt{2} < 3$.

b) Hàm số nghịch biến khi $x > 0$ nên $f\left(2\sqrt{2}\right) > f(3)$ (kiểm tra: $f\left(2\sqrt{2}\right) = -\dfrac{8}{3}$, $f(3) = -3$).

c) Điểm có tung độ $-3$ thuộc đồ thị khi hoành độ $x$ thoả mãn $-\dfrac{1}{3}x^2 = -3 \Leftrightarrow x^2 = 9 \Leftrightarrow x = 3$ hoặc $x = -3$.

c) Vậy có hai điểm $(3; -3)$ và $(-3; -3)$; hai điểm này đối xứng nhau qua trục $Oy$, phù hợp với tính chất của đồ thị.

Đáp số: a) $\boxed{\text{Đồng biến khi } x < 0, \text{ nghịch biến khi } x > 0; \text{ parabol đỉnh } O, \text{ trục đối xứng } Oy, \text{ nằm phía dưới trục hoành}}$; b) $\boxed{f\left(2\sqrt{2}\right) > f(3)}$; c) $\boxed{(3; -3) \text{ và } (-3; -3)}$
\end{solution}
\end{problem}

% ---------------------------------------------------------------- chứng minh tính chất với tham số, hai cách
\begin{problem}[9DS-HSB2-004]
% nguon: dummy (tự sinh để test)
% lop: 9
% kc_du_kien: 9DS-HSB2-02, 8DS-HDT-01
% do_kho_du_kien: 3
Cho hàm số $y = (m^2 - 2m + 3)x^2$, với $m$ là tham số. Chứng minh rằng với mọi giá trị của $m$, hàm số luôn nghịch biến khi $x < 0$ và đồng biến khi $x > 0$.

\begin{solution}[cách 1: dùng tính chất của hàm số $y = ax^2$]
Hệ số $a = m^2 - 2m + 3$; tách và dùng hằng đẳng thức bình phương của một hiệu: $a = (m^2 - 2m + 1) + 2 = (m - 1)^2 + 2$.

Vì $(m - 1)^2 \geq 0$ với mọi $m$ nên $a \geq 2 > 0$ với mọi $m$; đặc biệt $a \neq 0$, nên hàm số có dạng $y = ax^2$ với $a > 0$.

Theo tính chất của hàm số $y = ax^2$ với $a > 0$: hàm số nghịch biến khi $x < 0$ và đồng biến khi $x > 0$.

Đáp số: $\boxed{a = (m - 1)^2 + 2 > 0 \text{ với mọi } m \text{ nên hàm số nghịch biến khi } x < 0, \text{ đồng biến khi } x > 0}$
\end{solution}

\begin{solution}[cách 2: dùng định nghĩa hàm số đồng biến, nghịch biến]
Như cách 1, $a = (m - 1)^2 + 2 > 0$ với mọi $m$. Đặt $f(x) = ax^2$.

Lấy hai giá trị bất kỳ $x_1 < x_2$ và xét hiệu (dùng hằng đẳng thức hiệu hai bình phương): $f(x_1) - f(x_2) = a(x_1^2 - x_2^2) = a(x_1 - x_2)(x_1 + x_2)$.

Nếu $x_1 < x_2 < 0$ thì $x_1 - x_2 < 0$ và $x_1 + x_2 < 0$, nên $(x_1 - x_2)(x_1 + x_2) > 0$; nhân với $a > 0$ được $f(x_1) - f(x_2) > 0$, tức là $f(x_1) > f(x_2)$. Vậy hàm số nghịch biến khi $x < 0$.

Nếu $0 < x_1 < x_2$ thì $x_1 - x_2 < 0$ và $x_1 + x_2 > 0$, nên $(x_1 - x_2)(x_1 + x_2) < 0$; nhân với $a > 0$ được $f(x_1) < f(x_2)$. Vậy hàm số đồng biến khi $x > 0$.

Đáp số: $\boxed{a = (m - 1)^2 + 2 > 0 \text{ với mọi } m \text{ nên hàm số nghịch biến khi } x < 0, \text{ đồng biến khi } x > 0}$
\end{solution}
\end{problem}

% ---------------------------------------------------------------- giao điểm; đường thẳng song song và tiếp xúc
\begin{problem}[9DS-HSB2-005]
% nguon: dummy (tự sinh để test)
% lop: 9
% kc_du_kien: 9DS-HSB2-03, 9DS-HSB2-04, 8DS-HSBN-04
% do_kho_du_kien: 3
Cho parabol $(P)\colon y = x^2$ và đường thẳng $(d)\colon y = x + 2$.

a) Tìm toạ độ các giao điểm của $(P)$ và $(d)$.

b) Viết phương trình đường thẳng $(d')$ song song với $(d)$ và tiếp xúc với $(P)$; tìm toạ độ tiếp điểm.

\begin{solution}
a) Hoành độ giao điểm của $(P)$ và $(d)$ là nghiệm của phương trình $x^2 = x + 2 \Leftrightarrow x^2 - x - 2 = 0$.

a) Phương trình có $a - b + c = 1 - (-1) + (-2) = 0$ nên có hai nghiệm $x_1 = -1$ và $x_2 = -\dfrac{c}{a} = 2$.

a) Tính tung độ từ $y = x^2$: với $x = -1$ thì $y = 1$; với $x = 2$ thì $y = 4$. Vậy hai giao điểm là $A(-1; 1)$ và $B(2; 4)$ (thử lại bằng $y = x + 2$: $-1 + 2 = 1$, $2 + 2 = 4$, đúng).

b) Đường thẳng $(d')$ song song với $(d)$ nên có cùng hệ số góc $1$ và tung độ gốc khác $2$: $(d')\colon y = x + n$ với $n \neq 2$.

b) Phương trình hoành độ giao điểm của $(P)$ và $(d')$: $x^2 = x + n \Leftrightarrow x^2 - x - n = 0$, có $\Delta = (-1)^2 - 4 \cdot 1 \cdot (-n) = 1 + 4n$.

b) $(d')$ tiếp xúc với $(P)$ khi và chỉ khi phương trình hoành độ có nghiệm kép, tức là $\Delta = 0 \Leftrightarrow 1 + 4n = 0 \Leftrightarrow n = -\dfrac{1}{4}$ (thoả mãn $n \neq 2$). Vậy $(d')\colon y = x - \dfrac{1}{4}$.

b) Khi đó nghiệm kép là $x = \dfrac{-(-1)}{2 \cdot 1} = \dfrac{1}{2}$, tung độ $y = \left(\dfrac{1}{2}\right)^2 = \dfrac{1}{4}$. Tiếp điểm là $\left(\dfrac{1}{2}; \dfrac{1}{4}\right)$.

Đáp số: a) $\boxed{A(-1; 1),\ B(2; 4)}$; b) $\boxed{(d')\colon y = x - \dfrac{1}{4},\ \text{tiếp điểm } \left(\dfrac{1}{2}; \dfrac{1}{4}\right)}$
\end{solution}
\end{problem}

% ---------------------------------------------------------------- tìm a rồi tìm giao điểm
\begin{problem}[9DS-HSB2-006]
% nguon: dummy (tự sinh để test)
% lop: 9
% kc_du_kien: 9DS-HSB2-01, 9DS-HSB2-03, 9DS-HSB2-02
% do_kho_du_kien: 3
Cho hàm số $y = ax^2$ có đồ thị là parabol $(P)$ đi qua điểm $M(2; -2)$.

a) Tìm $a$ và cho biết $(P)$ nằm phía trên hay phía dưới trục hoành.

b) Tìm toạ độ các giao điểm của $(P)$ với đường thẳng $(d)\colon y = x - 4$.

\begin{solution}
a) $M(2; -2)$ thuộc $(P)$ nên $-2 = a \cdot 2^2 = 4a$, suy ra $a = -\dfrac{1}{2}$ (thoả mãn $a \neq 0$); $(P)\colon y = -\dfrac{1}{2}x^2$.

a) Vì $a = -\dfrac{1}{2} < 0$ nên $(P)$ nằm phía dưới trục hoành, $O$ là điểm cao nhất của $(P)$.

b) Phương trình hoành độ giao điểm: $-\dfrac{1}{2}x^2 = x - 4$; nhân hai vế với $-2$ rồi chuyển vế: $x^2 + 2x - 8 = 0$.

b) Dùng công thức nghiệm thu gọn với $b' = 1$: $\Delta' = 1^2 - 1 \cdot (-8) = 9$, $\sqrt{\Delta'} = 3$, nên $x_1 = -1 + 3 = 2$ và $x_2 = -1 - 3 = -4$.

b) Tính tung độ từ $y = x - 4$: với $x = 2$ thì $y = -2$; với $x = -4$ thì $y = -8$.

b) Kiểm tra trên $(P)$: $-\dfrac{1}{2} \cdot 2^2 = -2$ và $-\dfrac{1}{2} \cdot (-4)^2 = -8$ (đúng). Giao điểm thứ nhất chính là $M(2; -2)$.

Đáp số: a) $\boxed{a = -\dfrac{1}{2},\ (P) \text{ nằm phía dưới trục hoành}}$; b) $\boxed{(2; -2) \text{ và } (-4; -8)}$
\end{solution}
\end{problem}

% ---------------------------------------------------------------- vị trí theo tham số, dấu các nghiệm
\begin{problem}[9DS-HSB2-007]
% nguon: dummy (tự sinh để test)
% lop: 9
% kc_du_kien: 9DS-HSB2-04, 9DS-PT2-03
% do_kho_du_kien: 4
Cho parabol $(P)\colon y = x^2$ và đường thẳng $(d)\colon y = 2x + m$, với $m$ là tham số.

a) Tìm $m$ để $(d)$ và $(P)$ không có điểm chung.

b) Tìm $m$ để $(d)$ cắt $(P)$ tại hai điểm phân biệt cùng nằm bên phải trục tung.

\begin{solution}
a) Phương trình hoành độ giao điểm của $(P)$ và $(d)$: $x^2 = 2x + m \Leftrightarrow x^2 - 2x - m = 0$ $(*)$. Số điểm chung của $(d)$ và $(P)$ bằng số nghiệm của $(*)$.

a) Phương trình $(*)$ có $b' = -1$, nên $\Delta' = (-1)^2 - 1 \cdot (-m) = 1 + m$.

a) $(d)$ và $(P)$ không có điểm chung khi và chỉ khi $(*)$ vô nghiệm, tức là $\Delta' < 0 \Leftrightarrow m < -1$.

b) Hai giao điểm nằm bên phải trục tung khi và chỉ khi $(*)$ có hai nghiệm phân biệt $x_1, x_2$ đều dương (hoành độ dương).

b) Điều kiện có hai nghiệm phân biệt: $\Delta' > 0 \Leftrightarrow m > -1$.

b) Khi đó, theo định lý Viète: $x_1 + x_2 = 2$ và $x_1 x_2 = -m$. Hai nghiệm cùng dương khi và chỉ khi $x_1 + x_2 > 0$ (luôn đúng vì $2 > 0$) và $x_1 x_2 > 0 \Leftrightarrow -m > 0 \Leftrightarrow m < 0$.

b) Kết hợp các điều kiện $m > -1$ và $m < 0$: $-1 < m < 0$.

Đáp số: a) $\boxed{m < -1}$; b) $\boxed{-1 < m < 0}$
\end{solution}
\end{problem}

% ---------------------------------------------------------------- luôn cắt, hệ thức Viète
\begin{problem}[9DS-HSB2-008]
% nguon: dummy (tự sinh để test)
% lop: 9
% kc_du_kien: 9DS-HSB2-04, 9DS-PT2-03
% do_kho_du_kien: 4
Cho parabol $(P)\colon y = x^2$ và đường thẳng $(d)\colon y = 2mx - 2m + 3$, với $m$ là tham số.

a) Chứng minh rằng $(d)$ luôn cắt $(P)$ tại hai điểm phân biệt với mọi $m$.

b) Gọi $x_1, x_2$ là hoành độ các giao điểm của $(d)$ và $(P)$. Tìm $m$ để $x_1^2 + x_2^2 = 14$.

\begin{solution}
a) Phương trình hoành độ giao điểm: $x^2 = 2mx - 2m + 3 \Leftrightarrow x^2 - 2mx + 2m - 3 = 0$ $(*)$.

a) Phương trình $(*)$ có $b' = -m$: $\Delta' = m^2 - (2m - 3) = m^2 - 2m + 3 = (m - 1)^2 + 2$.

a) Vì $(m - 1)^2 \geq 0$ nên $\Delta' \geq 2 > 0$ với mọi $m$; do đó $(*)$ luôn có hai nghiệm phân biệt, tức là $(d)$ luôn cắt $(P)$ tại hai điểm phân biệt.

b) Theo định lý Viète: $x_1 + x_2 = 2m$ và $x_1 x_2 = 2m - 3$.

b) Biến đổi theo tổng và tích: $x_1^2 + x_2^2 = (x_1 + x_2)^2 - 2x_1x_2 = 4m^2 - 2(2m - 3) = 4m^2 - 4m + 6$.

b) Điều kiện đề bài: $4m^2 - 4m + 6 = 14 \Leftrightarrow m^2 - m - 2 = 0$; phương trình này có $a - b + c = 1 + 1 - 2 = 0$ nên $m = -1$ hoặc $m = 2$.

b) Cả hai giá trị đều thoả mãn vì theo câu a) $(d)$ luôn cắt $(P)$ tại hai điểm phân biệt với mọi $m$. Thử lại với $m = 2$: $(*)$ là $x^2 - 4x + 1 = 0$, có $x_1^2 + x_2^2 = 16 - 2 = 14$ (đúng).

Đáp số: a) $\boxed{\Delta' = (m - 1)^2 + 2 > 0 \text{ với mọi } m}$; b) $\boxed{m = -1 \text{ hoặc } m = 2}$
\end{solution}
\end{problem}

% ---------------------------------------------------------------- tham số: hai phía trục tung, tung độ, diện tích nhỏ nhất
\begin{problem}[9DS-HSB2-009]
% nguon: dummy (tự sinh để test)
% lop: 9
% kc_du_kien: 9DS-HSB2-04, 9DS-HSB2-03, 9DS-PT2-03
% do_kho_du_kien: 5
Cho parabol $(P)\colon y = x^2$ và đường thẳng $(d)\colon y = mx + 1$, với $m$ là tham số.

a) Chứng minh rằng với mọi $m$, $(d)$ luôn cắt $(P)$ tại hai điểm phân biệt $A$, $B$ nằm về hai phía của trục tung.

b) Gọi $x_1, x_2$ lần lượt là hoành độ và $y_1, y_2$ lần lượt là tung độ của $A$, $B$. Tìm $m$ để $y_1 + y_2 = 3(x_1 + x_2)$.

c) Tìm $m$ để tam giác $OAB$ ($O$ là gốc toạ độ) có diện tích nhỏ nhất và tính diện tích nhỏ nhất đó.

\begin{solution}
a) Phương trình hoành độ giao điểm của $(P)$ và $(d)$: $x^2 = mx + 1 \Leftrightarrow x^2 - mx - 1 = 0$ $(*)$.

a) Phương trình $(*)$ có $\Delta = m^2 + 4 > 0$ với mọi $m$, nên luôn có hai nghiệm phân biệt $x_1, x_2$; tức là $(d)$ luôn cắt $(P)$ tại hai điểm phân biệt $A(x_1; y_1)$, $B(x_2; y_2)$.

a) Theo định lý Viète, $x_1 x_2 = -1 < 0$ nên $x_1$, $x_2$ trái dấu; vậy $A$ và $B$ nằm về hai phía của trục tung.

b) Theo định lý Viète: $x_1 + x_2 = m$.

b) $A$, $B$ thuộc $(d)$ nên $y_1 = mx_1 + 1$ và $y_2 = mx_2 + 1$; suy ra $y_1 + y_2 = m(x_1 + x_2) + 2 = m^2 + 2$.

b) Điều kiện đề bài: $m^2 + 2 = 3m \Leftrightarrow m^2 - 3m + 2 = 0$; vì $a + b + c = 1 - 3 + 2 = 0$ nên $m = 1$ hoặc $m = 2$. Cả hai giá trị đều thoả mãn vì câu a) đúng với mọi $m$.

c) $(d)$ cắt trục tung tại $C(0; 1)$, nên $OC = 1$. Vì $A$, $B$ nằm về hai phía của trục tung nên $C$ nằm giữa $A$ và $B$, và tam giác $OAB$ được chia thành hai tam giác $OAC$, $OBC$ có chung cạnh $OC$.

c) Khoảng cách từ $A$, $B$ đến trục tung (đường cao hạ xuống cạnh $OC$) lần lượt là $|x_1|$, $|x_2|$, nên $S_{OAB} = S_{OAC} + S_{OBC} = \dfrac{1}{2} \cdot OC \cdot (|x_1| + |x_2|) = \dfrac{1}{2}(|x_1| + |x_2|)$.

c) Vì $x_1$, $x_2$ trái dấu nên $|x_1| + |x_2| = |x_1 - x_2|$; mà $(x_1 - x_2)^2 = (x_1 + x_2)^2 - 4x_1x_2 = m^2 + 4$, nên $|x_1 - x_2| = \sqrt{m^2 + 4}$ và $S_{OAB} = \dfrac{1}{2}\sqrt{m^2 + 4}$.

c) Vì $m^2 \geq 0$ nên $\sqrt{m^2 + 4} \geq \sqrt{4} = 2$, suy ra $S_{OAB} \geq 1$; dấu bằng xảy ra khi và chỉ khi $m = 0$.

Đáp số: a) $\boxed{\Delta = m^2 + 4 > 0,\ x_1 x_2 = -1 < 0}$; b) $\boxed{m = 1 \text{ hoặc } m = 2}$; c) $\boxed{m = 0,\ S_{\min} = 1 \text{ (đơn vị diện tích)}}$
\end{solution}
\end{problem}

% ---------------------------------------------------------------- bài toán thực tế: lực gió
\begin{problem}[9DS-HSB2-010]
% nguon: dummy (tự sinh để test)
% lop: 9
% kc_du_kien: 9DS-HSB2-01
% do_kho_du_kien: 2
Lực $F$ (N) của gió tác động lên một cánh buồm tỉ lệ thuận với bình phương vận tốc $v$ (m/s) của gió, tức là $F = av^2$ ($a$ là hằng số). Biết khi vận tốc gió bằng $2$ m/s thì lực tác động lên cánh buồm bằng $120$ N.

a) Tìm hằng số $a$.

b) Tính lực $F$ khi vận tốc gió bằng $10$ m/s.

c) Cánh buồm chỉ chịu được lực tối đa $12\,000$ N. Hỏi thuyền có thể đi được trong gió bão với vận tốc $90$ km/h hay không?

\begin{solution}
a) Thay $v = 2$, $F = 120$ vào công thức $F = av^2$: $120 = a \cdot 2^2 = 4a$, suy ra $a = 30$.

b) Với $a = 30$ và $v = 10$: $F = 30 \cdot 10^2 = 3\,000$ (N).

c) Đổi đơn vị vận tốc: $90$ km/h $= \dfrac{90\,000 \text{ m}}{3\,600 \text{ s}} = 25$ m/s.

c) Lực của gió khi đó: $F = 30 \cdot 25^2 = 30 \cdot 625 = 18\,750$ (N).

c) Vì $18\,750 > 12\,000$ nên cánh buồm không chịu được lực này; thuyền không thể đi trong gió bão $90$ km/h. (Cánh buồm chỉ chịu được khi $30v^2 \leq 12\,000$, tức là $v^2 \leq 400$, hay $v \leq 20$ m/s.)

Đáp số: a) $\boxed{a = 30}$; b) $\boxed{F = 3\,000 \text{ N}}$; c) $\boxed{\text{Không, vì } F = 18\,750 \text{ N} > 12\,000 \text{ N}}$
\end{solution}
\end{problem}

% ---------------------------------------------------------------- bảng giá trị, tính đối xứng của đồ thị
\begin{problem}[9DS-HSB2-011]
% nguon: dummy (tự sinh để test)
% lop: 9
% kc_du_kien: 9DS-HSB2-01, 9DS-HSB2-02
% do_kho_du_kien: 1
Cho hàm số $y = -\dfrac{1}{2}x^2$.

a) Tính các giá trị tương ứng của $y$ khi $x$ lần lượt bằng $-4$; $-2$; $0$; $2$; $4$.

b) Nêu nhận xét về các cặp điểm của đồ thị có hoành độ đối nhau, và về vị trí của đồ thị so với trục hoành.

\begin{solution}
a) Thay từng giá trị của $x$ vào công thức $y = -\dfrac{1}{2}x^2$. Với $x = -4$ và $x = 4$: $(-4)^2 = 4^2 = 16$, nên $y = -\dfrac{1}{2} \cdot 16 = -8$.

a) Với $x = -2$ và $x = 2$: $(-2)^2 = 2^2 = 4$, nên $y = -\dfrac{1}{2} \cdot 4 = -2$. Với $x = 0$: $y = 0$.

b) Các điểm $(-4; -8)$ và $(4; -8)$, $(-2; -2)$ và $(2; -2)$ có hoành độ đối nhau và tung độ bằng nhau, nên đối xứng nhau qua trục $Oy$. Điều này đúng với mọi cặp hoành độ đối nhau vì $(-x)^2 = x^2$; do đó đồ thị nhận trục $Oy$ làm trục đối xứng.

b) Vì $a = -\dfrac{1}{2} < 0$ và $x^2 \geq 0$ nên $y = -\dfrac{1}{2}x^2 \leq 0$ với mọi $x$, dấu bằng chỉ khi $x = 0$; vậy đồ thị nằm phía dưới trục hoành và $O(0; 0)$ là điểm cao nhất của đồ thị.

Đáp số: a) $\boxed{y = -8;\ -2;\ 0;\ -2;\ -8}$; b) $\boxed{\text{Các điểm có hoành độ đối nhau đối xứng qua } Oy;\ \text{đồ thị nằm phía dưới trục hoành, } O \text{ là điểm cao nhất}}$
\end{solution}
\end{problem}

% ---------------------------------------------------------------- tham số trong hệ số a: vị trí đồ thị, đi qua điểm
\begin{problem}[9DS-HSB2-012]
% nguon: dummy (tự sinh để test)
% lop: 9
% kc_du_kien: 9DS-HSB2-02, 9DS-HSB2-01, 9DS-BPT-02
% do_kho_du_kien: 2
Cho hàm số $y = (2m - 3)x^2$, với $m$ là tham số, $m \neq \dfrac{3}{2}$.

a) Tìm $m$ để đồ thị của hàm số nằm phía trên trục hoành.

b) Tìm $m$ để đồ thị của hàm số đi qua điểm $A(-2; -6)$.

\begin{solution}
a) Hàm số có dạng $y = ax^2$ với $a = 2m - 3 \neq 0$. Đồ thị của hàm số $y = ax^2$ nằm phía trên trục hoành khi và chỉ khi $a > 0$.

a) Giải bất phương trình $2m - 3 > 0$: chuyển vế được $2m > 3$, chia hai vế cho $2 > 0$ (giữ nguyên chiều) được $m > \dfrac{3}{2}$.

b) Đồ thị đi qua $A(-2; -6)$ khi và chỉ khi toạ độ của $A$ thoả mãn công thức hàm số: $-6 = (2m - 3) \cdot (-2)^2 = 4(2m - 3)$.

b) Chia hai vế cho $4$: $2m - 3 = -\dfrac{3}{2}$, suy ra $2m = \dfrac{3}{2}$, $m = \dfrac{3}{4}$ (thoả mãn $m \neq \dfrac{3}{2}$). Khi đó $a = -\dfrac{3}{2} < 0$, đồ thị nằm phía dưới trục hoành, phù hợp với việc $A$ có tung độ âm.

Đáp số: a) $\boxed{m > \dfrac{3}{2}}$; b) $\boxed{m = \dfrac{3}{4}}$
\end{solution}
\end{problem}

% ---------------------------------------------------------------- bài toán thực tế: vật rơi tự do
\begin{problem}[9DS-HSB2-013]
% nguon: dummy (tự sinh để test)
% lop: 9
% kc_du_kien: 9DS-HSB2-01, 9DS-PT2-01
% do_kho_du_kien: 2
Khi thả rơi một vật từ trên cao (bỏ qua sức cản của không khí), quãng đường $s$ (m) vật rơi được sau $t$ giây được cho bởi công thức $s = 5t^2$.

a) Sau $3$ giây, vật rơi được bao nhiêu mét?

b) Một vật được thả rơi từ đỉnh một toà tháp cao $125$ m. Sau bao lâu thì vật chạm đất?

c) Sau $2$ giây kể từ lúc thả, vật ở câu b) còn cách mặt đất bao nhiêu mét?

\begin{solution}
a) Thay $t = 3$ vào công thức: $s = 5 \cdot 3^2 = 5 \cdot 9 = 45$ (m).

b) Gọi $t$ (giây, $t > 0$) là thời gian từ lúc thả đến lúc vật chạm đất; khi đó vật rơi được đúng $125$ m, nên $5t^2 = 125$.

b) Chia hai vế cho $5$: $t^2 = 25$, suy ra $t = 5$ hoặc $t = -5$. Vì $t > 0$ nên $t = 5$ (giây).

c) Sau $2$ giây vật rơi được $s = 5 \cdot 2^2 = 20$ (m); vì $2 < 5$ nên lúc đó vật chưa chạm đất.

c) Khoảng cách từ vật đến mặt đất bằng độ cao ban đầu trừ quãng đường đã rơi: $125 - 20 = 105$ (m).

Đáp số: a) $\boxed{45 \text{ m}}$; b) $\boxed{5 \text{ giây}}$; c) $\boxed{105 \text{ m}}$
\end{solution}
\end{problem}

% ---------------------------------------------------------------- phương trình đường thẳng qua hai điểm của parabol, hai cách
\begin{problem}[9DS-HSB2-014]
% nguon: dummy (tự sinh để test)
% lop: 9
% kc_du_kien: 9DS-HSB2-01, 9DS-HPT-08, 9DS-PT2-03
% do_kho_du_kien: 3
Cho parabol $(P)\colon y = \dfrac{1}{2}x^2$. Gọi $A$, $B$ là hai điểm thuộc $(P)$ có hoành độ lần lượt là $-2$ và $4$.

a) Tìm toạ độ của $A$ và $B$.

b) Viết phương trình đường thẳng $AB$.

\begin{solution}[cách 1: lập hệ phương trình]
a) $A$ thuộc $(P)$ và có hoành độ $-2$ nên tung độ của $A$ là $y = \dfrac{1}{2} \cdot (-2)^2 = 2$; vậy $A(-2; 2)$.

a) $B$ thuộc $(P)$ và có hoành độ $4$ nên tung độ của $B$ là $y = \dfrac{1}{2} \cdot 4^2 = 8$; vậy $B(4; 8)$.

b) $A$ và $B$ có hoành độ khác nhau nên đường thẳng $AB$ không song song với trục tung, có dạng $y = ax + b$. Vì $A$, $B$ thuộc đường thẳng này nên $2 = -2a + b$ $(1)$ và $8 = 4a + b$ $(2)$.

b) Giải hệ bằng phương pháp cộng đại số: trừ vế theo vế $(2)$ cho $(1)$ được $6 = 6a$, suy ra $a = 1$; thay vào $(1)$: $b = 2 + 2a = 4$.

b) Vậy $AB: y = x + 4$. Kiểm tra: $-2 + 4 = 2$ và $4 + 4 = 8$ (đúng).

Đáp số: a) $\boxed{A(-2; 2),\ B(4; 8)}$; b) $\boxed{y = x + 4}$
\end{solution}

\begin{solution}[cách 2: dùng định lý Viète]
a) Thay hoành độ vào công thức của $(P)$: $y_A = \dfrac{1}{2} \cdot (-2)^2 = 2$ và $y_B = \dfrac{1}{2} \cdot 4^2 = 8$; vậy $A(-2; 2)$, $B(4; 8)$.

b) Đường thẳng $AB$ không song song với trục tung (vì $A$, $B$ có hoành độ khác nhau) nên có dạng $y = ax + b$. Hoành độ giao điểm của $AB$ và $(P)$ là nghiệm của phương trình $\dfrac{1}{2}x^2 = ax + b \Leftrightarrow x^2 - 2ax - 2b = 0$ $(*)$.

b) $A$ và $B$ là hai điểm chung của $AB$ và $(P)$, nên $-2$ và $4$ là nghiệm của $(*)$; phương trình bậc hai $(*)$ có không quá hai nghiệm, nên đó là hai nghiệm của $(*)$. Theo định lý Viète: $-2 + 4 = 2a$ và $(-2) \cdot 4 = -2b$.

b) Suy ra $a = 1$ và $b = 4$; vậy $AB: y = x + 4$.

Đáp số: a) $\boxed{A(-2; 2),\ B(4; 8)}$; b) $\boxed{y = x + 4}$
\end{solution}
\end{problem}

% ---------------------------------------------------------------- đường thẳng y = m, tam giác vuông tại O, hai cách
\begin{problem}[9DS-HSB2-015]
% nguon: dummy (tự sinh để test)
% lop: 9
% kc_du_kien: 9DS-HSB2-02, 9DS-HSB2-03
% do_kho_du_kien: 3
Cho parabol $(P)\colon y = x^2$ và đường thẳng $(d)\colon y = m$ (với $m > 0$).

a) Chứng minh rằng $(d)$ luôn cắt $(P)$ tại hai điểm phân biệt $A$, $B$ đối xứng nhau qua trục $Oy$.

b) Tìm $m$ để tam giác $OAB$ vuông tại $O$ ($O$ là gốc toạ độ). Khi đó tính diện tích tam giác $OAB$.

\begin{solution}[cách 1: dùng đường trung tuyến]
a) Hoành độ giao điểm của $(P)$ và $(d)$ là nghiệm của phương trình $x^2 = m$. Vì $m > 0$ nên phương trình có hai nghiệm phân biệt $x = -\sqrt{m}$ và $x = \sqrt{m}$; hai giao điểm là $A\left(-\sqrt{m}; m\right)$ và $B\left(\sqrt{m}; m\right)$.

a) $A$ và $B$ có hoành độ đối nhau và tung độ bằng nhau nên đối xứng nhau qua trục $Oy$ (phù hợp với tính chất đồ thị hàm số $y = ax^2$ nhận $Oy$ làm trục đối xứng).

b) Gọi $H(0; m)$ là giao điểm của $(d)$ với trục tung. Vì $A$, $B$ đối xứng nhau qua $Oy$ nên $H$ là trung điểm của $AB$, với $HA = HB = \sqrt{m}$ và $OH = m$; do đó $OH$ là đường trung tuyến của tam giác $OAB$.

b) Tam giác có đường trung tuyến ứng với một cạnh bằng nửa cạnh ấy thì là tam giác vuông (và ngược lại), nên tam giác $OAB$ vuông tại $O$ khi và chỉ khi $OH = \dfrac{AB}{2} = HA$, tức là $m = \sqrt{m}$.

b) Vì $m > 0$, chia hai vế cho $\sqrt{m} > 0$ được $\sqrt{m} = 1$, suy ra $m = 1$.

b) Khi $m = 1$: $A(-1; 1)$, $B(1; 1)$, $AB = 2$, $OH = 1$ và $OH \perp AB$ (vì $AB$ nằm trên đường thẳng $y = 1$ song song với trục hoành). Diện tích $S_{OAB} = \dfrac{1}{2} \cdot AB \cdot OH = \dfrac{1}{2} \cdot 2 \cdot 1 = 1$.

Đáp số: a) $\boxed{A\left(-\sqrt{m}; m\right),\ B\left(\sqrt{m}; m\right) \text{ đối xứng nhau qua } Oy}$; b) $\boxed{m = 1,\ S_{OAB} = 1 \text{ (đơn vị diện tích)}}$
\end{solution}

\begin{solution}[cách 2: dùng định lí Pythagore]
a) Phương trình hoành độ giao điểm $x^2 = m$ (với $m > 0$) có hai nghiệm phân biệt $x = \pm\sqrt{m}$, nên $(d)$ cắt $(P)$ tại $A\left(-\sqrt{m}; m\right)$ và $B\left(\sqrt{m}; m\right)$; hai điểm này có hoành độ đối nhau, tung độ bằng nhau nên đối xứng nhau qua $Oy$.

b) Gọi $H(0; m)$. Tam giác $OHA$ vuông tại $H$ nên theo định lí Pythagore $OA^2 = OH^2 + HA^2 = m^2 + m$; tương tự $OB^2 = m^2 + m$; và $AB = 2\sqrt{m}$ nên $AB^2 = 4m$.

b) Theo định lí Pythagore và định lí đảo, tam giác $OAB$ vuông tại $O$ khi và chỉ khi $OA^2 + OB^2 = AB^2$, tức là $2(m^2 + m) = 4m \Leftrightarrow 2m^2 - 2m = 0 \Leftrightarrow 2m(m - 1) = 0$.

b) Suy ra $m = 0$ (loại vì $m > 0$) hoặc $m = 1$ (thoả mãn). Khi đó $AB = 2$, đường cao $OH = 1$, nên $S_{OAB} = \dfrac{1}{2} \cdot 2 \cdot 1 = 1$.

Đáp số: a) $\boxed{A\left(-\sqrt{m}; m\right),\ B\left(\sqrt{m}; m\right) \text{ đối xứng nhau qua } Oy}$; b) $\boxed{m = 1,\ S_{OAB} = 1 \text{ (đơn vị diện tích)}}$
\end{solution}
\end{problem}

% ---------------------------------------------------------------- nhẩm nghiệm phương trình hoành độ, tiếp xúc, hoành độ gấp đôi
\begin{problem}[9DS-HSB2-016]
% nguon: dummy (tự sinh để test)
% lop: 9
% kc_du_kien: 9DS-HSB2-04, 9DS-HSB2-03, 9DS-PT2-04
% do_kho_du_kien: 3
Cho parabol $(P)\colon y = x^2$ và đường thẳng $(d)\colon y = (m + 1)x - m$, với $m$ là tham số.

a) Chứng minh rằng với mọi $m$, $(d)$ và $(P)$ luôn có điểm chung $A(1; 1)$. Tìm $m$ để $(d)$ tiếp xúc với $(P)$.

b) Tìm $m$ để $(d)$ cắt $(P)$ tại hai điểm phân biệt mà hoành độ của điểm này gấp đôi hoành độ của điểm kia.

\begin{solution}
a) Phương trình hoành độ giao điểm của $(P)$ và $(d)$: $x^2 = (m + 1)x - m \Leftrightarrow x^2 - (m + 1)x + m = 0$ $(*)$.

a) Phương trình $(*)$ có $a + b + c = 1 - (m + 1) + m = 0$ nên luôn có hai nghiệm $x_1 = 1$ và $x_2 = \dfrac{c}{a} = m$. Với $x = 1$ thì $y = 1^2 = 1$, nên $(d)$ và $(P)$ luôn có điểm chung $A(1; 1)$ (kiểm tra trên $(d)$: $(m + 1) \cdot 1 - m = 1$).

a) $(d)$ tiếp xúc với $(P)$ khi và chỉ khi $(*)$ có nghiệm kép, tức là hai nghiệm $1$ và $m$ trùng nhau: $m = 1$ (khớp với $\Delta = (m + 1)^2 - 4m = (m - 1)^2 = 0$). Khi đó $(d)\colon y = 2x - 1$ và tiếp điểm là $A(1; 1)$.

b) $(d)$ cắt $(P)$ tại hai điểm phân biệt khi và chỉ khi $(*)$ có hai nghiệm phân biệt, tức là $m \neq 1$; khi đó hoành độ hai giao điểm là $1$ và $m$.

b) Hoành độ của điểm này gấp đôi hoành độ của điểm kia xảy ra trong hai trường hợp: $m = 2 \cdot 1 = 2$, hoặc $1 = 2m$, tức là $m = \dfrac{1}{2}$. Cả hai giá trị đều thoả mãn $m \neq 1$.

b) Thử lại: với $m = 2$, hai giao điểm là $(1; 1)$ và $(2; 4)$; với $m = \dfrac{1}{2}$, hai giao điểm là $(1; 1)$ và $\left(\dfrac{1}{2}; \dfrac{1}{4}\right)$ (đúng).

Đáp số: a) $\boxed{A(1; 1) \text{ luôn là điểm chung};\ \text{tiếp xúc khi } m = 1}$; b) $\boxed{m = 2 \text{ hoặc } m = \dfrac{1}{2}}$
\end{solution}
\end{problem}

% ---------------------------------------------------------------- đường thẳng qua một điểm cho trước tiếp xúc với parabol
\begin{problem}[9DS-HSB2-017]
% nguon: dummy (tự sinh để test)
% lop: 9
% kc_du_kien: 9DS-HSB2-04, 9DS-PT2-02, 8DS-HSBN-01
% do_kho_du_kien: 4
Cho parabol $(P)\colon y = x^2$ và điểm $A(1; -3)$. Viết phương trình các đường thẳng có dạng $y = ax + b$ đi qua $A$ và tiếp xúc với $(P)$; tìm toạ độ các tiếp điểm.

\begin{solution}
Đường thẳng $(d)\colon y = ax + b$ đi qua $A(1; -3)$ nên $-3 = a \cdot 1 + b$, suy ra $b = -a - 3$; vậy $(d)\colon y = ax - a - 3$.

Phương trình hoành độ giao điểm của $(d)$ và $(P)$: $x^2 = ax - a - 3 \Leftrightarrow x^2 - ax + a + 3 = 0$ $(*)$.

$(d)$ tiếp xúc với $(P)$ khi và chỉ khi $(*)$ có nghiệm kép, tức là $\Delta = (-a)^2 - 4 \cdot 1 \cdot (a + 3) = a^2 - 4a - 12 = 0$.

Giải phương trình bậc hai ẩn $a$ bằng công thức nghiệm thu gọn ($b' = -2$): $\Delta' = (-2)^2 - 1 \cdot (-12) = 16$, $\sqrt{\Delta'} = 4$, nên $a = 2 + 4 = 6$ hoặc $a = 2 - 4 = -2$.

Với $a = 6$: $b = -6 - 3 = -9$, ta được $(d_1)\colon y = 6x - 9$; nghiệm kép của $(*)$ là $x = \dfrac{a}{2} = 3$, tung độ $y = 3^2 = 9$, nên tiếp điểm là $(3; 9)$.

Với $a = -2$: $b = 2 - 3 = -1$, ta được $(d_2)\colon y = -2x - 1$; nghiệm kép là $x = \dfrac{a}{2} = -1$, tung độ $y = (-1)^2 = 1$, nên tiếp điểm là $(-1; 1)$.

Kiểm tra: với $(d_1)$, $(*)$ là $x^2 - 6x + 9 = (x - 3)^2 = 0$; với $(d_2)$, $(*)$ là $x^2 + 2x + 1 = (x + 1)^2 = 0$; cả hai đường thẳng đều đi qua $A$ vì $6 - 9 = -3$ và $-2 - 1 = -3$.

Đáp số: $\boxed{y = 6x - 9 \text{ (tiếp điểm } (3; 9)\text{)};\ y = -2x - 1 \text{ (tiếp điểm } (-1; 1)\text{)}}$
\end{solution}
\end{problem}

% ---------------------------------------------------------------- biểu thức không đối xứng của toạ độ giao điểm
\begin{problem}[9DS-HSB2-018]
% nguon: dummy (tự sinh để test)
% lop: 9
% kc_du_kien: 9DS-HSB2-04, 9DS-HSB2-03, 9DS-PT2-03
% do_kho_du_kien: 4
Cho parabol $(P)\colon y = x^2$ và đường thẳng $(d)\colon y = 2x - m + 3$, với $m$ là tham số.

a) Tìm $m$ để $(d)$ cắt $(P)$ tại hai điểm phân biệt.

b) Gọi $A(x_1; y_1)$ và $B(x_2; y_2)$ là hai giao điểm đó. Tìm $m$ để $y_1 + 2x_2 = 9$.

\begin{solution}
a) Phương trình hoành độ giao điểm của $(P)$ và $(d)$: $x^2 = 2x - m + 3 \Leftrightarrow x^2 - 2x + m - 3 = 0$ $(*)$; phương trình có $b' = -1$ nên $\Delta' = (-1)^2 - 1 \cdot (m - 3) = 4 - m$.

a) $(d)$ cắt $(P)$ tại hai điểm phân biệt khi và chỉ khi $(*)$ có hai nghiệm phân biệt, tức là $\Delta' > 0 \Leftrightarrow 4 - m > 0 \Leftrightarrow m < 4$.

b) Với $m < 4$, $x_1$, $x_2$ là hai nghiệm của $(*)$; theo định lý Viète: $x_1 + x_2 = 2$.

b) Điểm $A(x_1; y_1)$ thuộc $(d)$ nên $y_1 = 2x_1 - m + 3$ (cũng có thể thấy từ $y_1 = x_1^2$ và $x_1^2 = 2x_1 - m + 3$ vì $x_1$ là nghiệm của $(*)$).

b) Do đó $y_1 + 2x_2 = 2x_1 - m + 3 + 2x_2 = 2(x_1 + x_2) - m + 3 = 4 - m + 3 = 7 - m$.

b) Điều kiện đề bài: $7 - m = 9$, suy ra $m = -2$ (thoả mãn $m < 4$).

b) Thử lại: với $m = -2$, $(*)$ là $x^2 - 2x - 5 = 0$ có hai nghiệm $1 + \sqrt{6}$ và $1 - \sqrt{6}$. Nếu $x_1 = 1 + \sqrt{6}$, $x_2 = 1 - \sqrt{6}$ thì $y_1 + 2x_2 = 7 + 2\sqrt{6} + 2 - 2\sqrt{6} = 9$; nếu đổi vai trò, $y_1 + 2x_2 = 7 - 2\sqrt{6} + 2 + 2\sqrt{6} = 9$ (đúng).

Đáp số: a) $\boxed{m < 4}$; b) $\boxed{m = -2}$
\end{solution}
\end{problem}

% ---------------------------------------------------------------- bài toán thực tế: cổng hình parabol
\begin{problem}[9DS-HSB2-019]
% nguon: dummy (tự sinh để test)
% lop: 9
% kc_du_kien: 9DS-HSB2-01, 9DS-HSB2-02, 9DS-BPT-05
% do_kho_du_kien: 4
Một chiếc cổng có dạng parabol, cao $8$ m (tính từ đỉnh cổng đến mặt đất), khoảng cách giữa hai chân cổng là $8$ m. Biết rằng trong một hệ trục toạ độ $Oxy$ thích hợp (đơn vị trên các trục là mét), cổng là một phần của parabol $y = ax^2$.

a) Hãy chọn hệ trục toạ độ và tìm $a$.

b) Một xe tải rộng $6$ m, cao $4$ m đi vào chính giữa cổng. Hỏi xe có đi qua cổng được không?

c) Một xe tải khác cao $6$ m đi vào chính giữa cổng. Hỏi xe đó rộng tối đa bao nhiêu mét thì đi qua được cổng?

\begin{solution}
a) Chọn hệ trục $Oxy$ có gốc $O$ tại đỉnh cổng, trục $Oy$ là trục đối xứng của cổng, hướng lên trên, trục $Ox$ nằm ngang. Khi đó cổng là một phần của parabol $y = ax^2$ (đỉnh $O$, nhận $Oy$ làm trục đối xứng), nằm phía dưới đỉnh nên $a < 0$; mặt đất là đường thẳng $y = -8$.

a) Hai chân cổng cách nhau $8$ m và đối xứng nhau qua $Oy$ nên có hoành độ $-4$ và $4$; chân cổng bên phải là điểm $(4; -8)$ thuộc parabol, nên $-8 = a \cdot 4^2 = 16a$, suy ra $a = -\dfrac{1}{2}$ (thoả mãn $a < 0$).

b) Xe rộng $6$ m đi chính giữa cổng nên hai mép thùng xe ở các hoành độ $x = -3$ và $x = 3$; do tính đối xứng của parabol qua $Oy$, chỉ cần xét $x = 3$. Tại $x = 3$, điểm của cổng có tung độ $y = -\dfrac{1}{2} \cdot 3^2 = -\dfrac{9}{2}$, tức là cách mặt đất $-\dfrac{9}{2} - (-8) = \dfrac{7}{2}$ (m).

b) Xe cao $4$ m mà tại mép thùng xe cổng chỉ cao $\dfrac{7}{2}$ m $< 4$ m, nên mép trên của thùng xe sẽ chạm vào cổng; vậy xe không đi qua được.

c) Gọi chiều rộng của xe là $w$ (m, $w > 0$); hai mép thùng xe ở các hoành độ $x = \pm\dfrac{w}{2}$. Với $a < 0$, khi $|x|$ tăng thì $y = -\dfrac{1}{2}x^2$ giảm, nên phần cổng phía trên thùng xe thấp nhất tại hai mép. Xe đi qua được khi tại $x = \dfrac{w}{2}$ độ cao của cổng so với mặt đất không nhỏ hơn $6$ m: $-\dfrac{1}{2} \cdot \left(\dfrac{w}{2}\right)^2 + 8 \geq 6$.

c) Biến đổi: $\dfrac{1}{2} \cdot \dfrac{w^2}{4} \leq 2 \Leftrightarrow w^2 \leq 16$. Vì $w > 0$ và $16 = 4^2$ nên $w^2 \leq 16 \Leftrightarrow w \leq 4$. Vậy xe rộng tối đa $4$ m (khi $w = 4$, tại mép xe cổng cao đúng $-\dfrac{1}{2} \cdot 2^2 + 8 = 6$ m).

Đáp số: a) $\boxed{a = -\dfrac{1}{2}}$; b) $\boxed{\text{Không đi qua được (tại mép xe cổng chỉ cao } \dfrac{7}{2} \text{ m} < 4 \text{ m)}}$; c) $\boxed{4 \text{ m}}$
\end{solution}
\end{problem}

% ---------------------------------------------------------------- độ dài dây cung qua điểm cố định, chứng minh hệ thức
\begin{problem}[9DS-HSB2-020]
% nguon: dummy (tự sinh để test)
% lop: 9
% kc_du_kien: 9DS-HSB2-04, 9DS-HSB2-03, 9DS-PT2-03
% do_kho_du_kien: 5
Cho parabol $(P)\colon y = \dfrac{1}{4}x^2$ và đường thẳng $(d)\colon y = mx + 1$, với $m$ là tham số.

a) Chứng minh rằng với mọi $m$, $(d)$ luôn cắt $(P)$ tại hai điểm phân biệt $A(x_1; y_1)$ và $B(x_2; y_2)$.

b) Chứng minh rằng $AB = 4(m^2 + 1)$. Tìm $m$ để $AB = 8$.

c) Chứng minh rằng $AB = y_1 + y_2 + 2$.

\begin{solution}
a) Phương trình hoành độ giao điểm của $(P)$ và $(d)$: $\dfrac{1}{4}x^2 = mx + 1$; nhân hai vế với $4$ rồi chuyển vế: $x^2 - 4mx - 4 = 0$ $(*)$.

a) Phương trình $(*)$ có $b' = -2m$ nên $\Delta' = (-2m)^2 - 1 \cdot (-4) = 4m^2 + 4 > 0$ với mọi $m$; do đó $(*)$ luôn có hai nghiệm phân biệt $x_1$, $x_2$, tức là $(d)$ luôn cắt $(P)$ tại hai điểm phân biệt $A(x_1; y_1)$, $B(x_2; y_2)$.

b) Theo định lý Viète: $x_1 + x_2 = 4m$ và $x_1x_2 = -4$. Vì $A$, $B$ thuộc $(d)$ nên $y_1 = mx_1 + 1$, $y_2 = mx_2 + 1$, suy ra $y_1 - y_2 = m(x_1 - x_2)$.

b) Gọi $K(x_2; y_1)$; khi đó $AK$ nằm trên đường thẳng song song với trục hoành và $BK$ nằm trên đường thẳng song song với trục tung, với $AK = |x_1 - x_2|$, $BK = |y_1 - y_2|$. Nếu $m \neq 0$ thì tam giác $AKB$ vuông tại $K$, theo định lí Pythagore $AB^2 = AK^2 + BK^2 = (x_1 - x_2)^2 + (y_1 - y_2)^2$; nếu $m = 0$ thì $y_1 = y_2$, $K$ trùng $B$ và đẳng thức này vẫn đúng.

b) Thay $y_1 - y_2 = m(x_1 - x_2)$: $AB^2 = (x_1 - x_2)^2(1 + m^2)$. Mặt khác $(x_1 - x_2)^2 = (x_1 + x_2)^2 - 4x_1x_2 = 16m^2 + 16 = 16(m^2 + 1)$, nên $AB^2 = 16(m^2 + 1)^2$; vì $AB > 0$ và $m^2 + 1 > 0$ nên $AB = 4(m^2 + 1)$.

b) $AB = 8 \Leftrightarrow 4(m^2 + 1) = 8 \Leftrightarrow m^2 = 1 \Leftrightarrow m = 1$ hoặc $m = -1$.

c) Từ $y_1 = mx_1 + 1$, $y_2 = mx_2 + 1$ và $x_1 + x_2 = 4m$: $y_1 + y_2 = m(x_1 + x_2) + 2 = 4m^2 + 2$.

c) Suy ra $y_1 + y_2 + 2 = 4m^2 + 4 = 4(m^2 + 1) = AB$ (theo câu b), điều phải chứng minh.

Đáp số: a) $\boxed{\Delta' = 4m^2 + 4 > 0 \text{ với mọi } m}$; b) $\boxed{AB = 4(m^2 + 1);\ m = 1 \text{ hoặc } m = -1}$; c) $\boxed{y_1 + y_2 + 2 = 4(m^2 + 1) = AB}$
\end{solution}
\end{problem}

% ---------------------------------------------------------------- nhận biết vị trí đồ thị, tính đồng biến qua dấu của a
\begin{problem}[9DS-HSB2-021]
% nguon: dummy (tự sinh để test)
% lop: 9
% kc_du_kien: 9DS-HSB2-02
% do_kho_du_kien: 1
Cho các hàm số: $y = 3x^2$; $y = -\dfrac{2}{5}x^2$; $y = \left(1 - \sqrt{2}\right)x^2$; $y = \dfrac{1}{4}x^2$.

a) Hàm số nào có đồ thị nằm phía trên trục hoành? Hàm số nào có đồ thị nằm phía dưới trục hoành?

b) Hàm số nào đồng biến khi $x > 0$? Hàm số nào nghịch biến khi $x > 0$?

\begin{solution}
a) Mỗi hàm số đều có dạng $y = ax^2$ với $a \neq 0$; theo tính chất của đồ thị, đồ thị nằm phía trên trục hoành khi $a > 0$ và nằm phía dưới trục hoành khi $a < 0$. Ta xét dấu của từng hệ số $a$.

a) Hàm số $y = 3x^2$ có $a = 3 > 0$; hàm số $y = \dfrac{1}{4}x^2$ có $a = \dfrac{1}{4} > 0$; hàm số $y = -\dfrac{2}{5}x^2$ có $a = -\dfrac{2}{5} < 0$.

a) Hàm số $y = \left(1 - \sqrt{2}\right)x^2$ có $a = 1 - \sqrt{2}$; vì $1 = \sqrt{1} < \sqrt{2}$ nên $a < 0$.

a) Vậy đồ thị của $y = 3x^2$ và $y = \dfrac{1}{4}x^2$ nằm phía trên trục hoành; đồ thị của $y = -\dfrac{2}{5}x^2$ và $y = \left(1 - \sqrt{2}\right)x^2$ nằm phía dưới trục hoành.

b) Theo tính chất của hàm số $y = ax^2$: nếu $a > 0$ thì hàm số đồng biến khi $x > 0$; nếu $a < 0$ thì hàm số nghịch biến khi $x > 0$. Dựa vào dấu của $a$ đã xét ở câu a), hai hàm số $y = 3x^2$, $y = \dfrac{1}{4}x^2$ đồng biến khi $x > 0$; hai hàm số $y = -\dfrac{2}{5}x^2$, $y = \left(1 - \sqrt{2}\right)x^2$ nghịch biến khi $x > 0$.

Đáp số: a) $\boxed{\text{Phía trên: } y = 3x^2,\ y = \dfrac{1}{4}x^2;\ \text{phía dưới: } y = -\dfrac{2}{5}x^2,\ y = \left(1 - \sqrt{2}\right)x^2}$; b) $\boxed{\text{Đồng biến khi } x > 0\colon y = 3x^2,\ y = \dfrac{1}{4}x^2;\ \text{nghịch biến khi } x > 0\colon y = -\dfrac{2}{5}x^2,\ y = \left(1 - \sqrt{2}\right)x^2}$
\end{solution}
\end{problem}

% ---------------------------------------------------------------- giao điểm với ba đường thẳng: cắt, tiếp xúc, không cắt
\begin{problem}[9DS-HSB2-022]
% nguon: dummy (tự sinh để test)
% lop: 9
% kc_du_kien: 9DS-HSB2-03, 9DS-HSB2-02, 9DS-PT2-04
% do_kho_du_kien: 2
Cho parabol $(P)\colon y = -\dfrac{1}{2}x^2$. Tìm toạ độ giao điểm (nếu có) của $(P)$ với mỗi đường thẳng sau:

a) $(d_1)\colon y = x - \dfrac{3}{2}$;

b) $(d_2)\colon y = 2x + 2$;

c) $(d_3)\colon y = 1$.

\begin{solution}
a) Hoành độ giao điểm của $(P)$ và $(d_1)$ là nghiệm của phương trình $-\dfrac{1}{2}x^2 = x - \dfrac{3}{2}$; nhân hai vế với $-2$ rồi chuyển vế: $x^2 + 2x - 3 = 0$.

a) Phương trình có $a + b + c = 1 + 2 - 3 = 0$ nên có hai nghiệm $x_1 = 1$ và $x_2 = \dfrac{c}{a} = -3$.

a) Tính tung độ từ $y = x - \dfrac{3}{2}$: với $x = 1$ thì $y = -\dfrac{1}{2}$; với $x = -3$ thì $y = -\dfrac{9}{2}$ (kiểm tra trên $(P)$: $-\dfrac{1}{2} \cdot 1^2 = -\dfrac{1}{2}$ và $-\dfrac{1}{2} \cdot (-3)^2 = -\dfrac{9}{2}$, đúng). Vậy $(d_1)$ cắt $(P)$ tại hai điểm $\left(1; -\dfrac{1}{2}\right)$ và $\left(-3; -\dfrac{9}{2}\right)$.

b) Phương trình hoành độ giao điểm của $(P)$ và $(d_2)$: $-\dfrac{1}{2}x^2 = 2x + 2$; nhân hai vế với $-2$ rồi chuyển vế: $x^2 + 4x + 4 = 0$, tức là $(x + 2)^2 = 0$ (hằng đẳng thức bình phương của một tổng), nên $x = -2$.

b) Phương trình có nghiệm kép $x = -2$, nên $(d_2)$ và $(P)$ có đúng một điểm chung ($(d_2)$ tiếp xúc với $(P)$). Tung độ: $y = 2 \cdot (-2) + 2 = -2$ (kiểm tra trên $(P)$: $-\dfrac{1}{2} \cdot (-2)^2 = -2$). Điểm chung là $(-2; -2)$.

c) Phương trình hoành độ giao điểm của $(P)$ và $(d_3)$: $-\dfrac{1}{2}x^2 = 1 \Leftrightarrow x^2 = -2$; phương trình này vô nghiệm vì $x^2 \geq 0$ với mọi $x$.

c) Vậy $(d_3)$ và $(P)$ không có điểm chung. Điều này phù hợp với tính chất của đồ thị: $a = -\dfrac{1}{2} < 0$ nên $(P)$ nằm phía dưới trục hoành, còn mọi điểm của $(d_3)$ có tung độ $1 > 0$.

Đáp số: a) $\boxed{\left(1; -\dfrac{1}{2}\right) \text{ và } \left(-3; -\dfrac{9}{2}\right)}$; b) $\boxed{(-2; -2)}$; c) $\boxed{\text{Không có giao điểm}}$
\end{solution}
\end{problem}

% ---------------------------------------------------------------- đọc bảng giá trị, rồi tìm giao điểm
\begin{problem}[9DS-HSB2-023]
% nguon: dummy (tự sinh để test)
% lop: 9
% kc_du_kien: 9DS-HSB2-03, 9DS-HSB2-01
% do_kho_du_kien: 2
Cho hàm số $y = ax^2$ ($a \neq 0$) có đồ thị là parabol $(P)$, và bảng giá trị chưa đầy đủ của hàm số như sau:
\[
\begin{array}{c|c|c|c|c|c}
x & -4 & -2 & 0 & 2 & 6 \\ \hline
y = ax^2 & \ldots & -1 & \ldots & \ldots & \ldots
\end{array}
\]

a) Tìm $a$ và điền các giá trị còn thiếu vào bảng.

b) Tìm toạ độ các giao điểm của $(P)$ với đường thẳng $(d)\colon y = -x - 3$. Các giao điểm đó có xuất hiện trong bảng giá trị ở câu a) không?

\begin{solution}
a) Theo bảng, khi $x = -2$ thì $y = -1$, tức là điểm $(-2; -1)$ thuộc $(P)$; thay vào công thức: $-1 = a \cdot (-2)^2 = 4a$, suy ra $a = -\dfrac{1}{4}$ (thoả mãn $a \neq 0$). Vậy $(P)\colon y = -\dfrac{1}{4}x^2$.

a) Thay lần lượt các giá trị còn lại của $x$ vào $y = -\dfrac{1}{4}x^2$: với $x = -4$ thì $y = -\dfrac{1}{4} \cdot 16 = -4$; với $x = 0$ thì $y = 0$; với $x = 2$ thì $y = -\dfrac{1}{4} \cdot 4 = -1$; với $x = 6$ thì $y = -\dfrac{1}{4} \cdot 36 = -9$.

b) Phương trình hoành độ giao điểm của $(P)$ và $(d)$: $-\dfrac{1}{4}x^2 = -x - 3$; nhân hai vế với $-4$ rồi chuyển vế: $x^2 - 4x - 12 = 0$.

b) Dùng công thức nghiệm thu gọn với $b' = -2$: $\Delta' = (-2)^2 - 1 \cdot (-12) = 16 > 0$, $\sqrt{\Delta'} = 4$, nên phương trình có hai nghiệm $x_1 = 2 + 4 = 6$ và $x_2 = 2 - 4 = -2$.

b) Tính tung độ từ $y = -x - 3$: với $x = 6$ thì $y = -9$; với $x = -2$ thì $y = -1$. Vậy hai giao điểm là $(6; -9)$ và $(-2; -1)$.

b) Đối chiếu với bảng ở câu a): cột $x = 6$ có $y = -9$ và cột $x = -2$ có $y = -1$, nên cả hai giao điểm đều xuất hiện trong bảng (đó là hai điểm của $(P)$ đồng thời thuộc $(d)$; các điểm còn lại của bảng không thuộc $(d)$, chẳng hạn với $x = 0$ thì $-x - 3 = -3 \neq 0$).

Đáp số: a) $\boxed{a = -\dfrac{1}{4};\ \text{các giá trị còn thiếu: } -4;\ 0;\ -1;\ -9}$; b) $\boxed{(6; -9) \text{ và } (-2; -1), \text{ cả hai đều có trong bảng}}$
\end{solution}
\end{problem}

% ---------------------------------------------------------------- biết một giao điểm, tìm a và giao điểm còn lại
\begin{problem}[9DS-HSB2-024]
% nguon: dummy (tự sinh để test)
% lop: 9
% kc_du_kien: 9DS-HSB2-03, 9DS-HSB2-01, 9DS-PT2-03
% do_kho_du_kien: 2
Đường thẳng $(d)\colon y = -x + 6$ cắt parabol $(P)\colon y = ax^2$ ($a \neq 0$) tại điểm $A$ có hoành độ bằng $-2$. Tìm $a$ và toạ độ giao điểm còn lại $B$ của $(d)$ và $(P)$.

\begin{solution}
$A$ thuộc $(d)$ và có hoành độ $-2$ nên tung độ của $A$ là $y = -(-2) + 6 = 8$; vậy $A(-2; 8)$.

$A$ cũng thuộc $(P)$ nên toạ độ của $A$ thoả mãn công thức của $(P)$: $8 = a \cdot (-2)^2 = 4a$, suy ra $a = 2$ (thoả mãn $a \neq 0$); $(P)\colon y = 2x^2$.

Phương trình hoành độ giao điểm của $(P)$ và $(d)$: $2x^2 = -x + 6 \Leftrightarrow 2x^2 + x - 6 = 0$ $(*)$.

Vì $A$ là một giao điểm nên $x_1 = -2$ là một nghiệm của $(*)$. Theo định lý Viète, $x_1 x_2 = \dfrac{c}{a} = \dfrac{-6}{2} = -3$, suy ra nghiệm còn lại là $x_2 = \dfrac{-3}{-2} = \dfrac{3}{2}$ (kiểm tra tổng: $-2 + \dfrac{3}{2} = -\dfrac{1}{2} = -\dfrac{b}{a}$, đúng).

Với $x = \dfrac{3}{2}$: $y = -\dfrac{3}{2} + 6 = \dfrac{9}{2}$ (kiểm tra trên $(P)$: $2 \cdot \left(\dfrac{3}{2}\right)^2 = 2 \cdot \dfrac{9}{4} = \dfrac{9}{2}$, đúng). Vậy $B\left(\dfrac{3}{2}; \dfrac{9}{2}\right)$.

Đáp số: $\boxed{a = 2,\ B\left(\dfrac{3}{2}; \dfrac{9}{2}\right)}$
\end{solution}
\end{problem}

% ---------------------------------------------------------------- giá trị lớn nhất, nhỏ nhất trên một đoạn
\begin{problem}[9DS-HSB2-025]
% nguon: dummy (tự sinh để test)
% lop: 9
% kc_du_kien: 9DS-HSB2-02, 9DS-HSB2-01
% do_kho_du_kien: 3
Cho hàm số $y = f(x) = -2x^2$.

a) Tìm giá trị lớn nhất và giá trị nhỏ nhất của hàm số khi $1 \leq x \leq 3$.

b) Tìm giá trị lớn nhất và giá trị nhỏ nhất của hàm số khi $-1 \leq x \leq 2$.

\begin{solution}
a) Hàm số có dạng $y = ax^2$ với $a = -2 < 0$, nên hàm số đồng biến khi $x < 0$ và nghịch biến khi $x > 0$.

a) Các giá trị $1 \leq x \leq 3$ đều dương, nên trên đoạn này hàm số nghịch biến: $x$ càng lớn thì $y$ càng nhỏ. Do đó với $1 \leq x \leq 3$ ta có $f(3) \leq f(x) \leq f(1)$, trong đó $f(1) = -2 \cdot 1^2 = -2$ và $f(3) = -2 \cdot 3^2 = -18$.

a) Vậy giá trị lớn nhất của hàm số là $-2$, đạt được khi $x = 1$; giá trị nhỏ nhất là $-18$, đạt được khi $x = 3$.

b) Đoạn $-1 \leq x \leq 2$ chứa $x = 0$, nên ta chia thành hai phần. Với $-1 \leq x \leq 0$, hàm số đồng biến nên $f(-1) \leq f(x) \leq f(0)$, tức là $-2 \leq f(x) \leq 0$.

b) Với $0 \leq x \leq 2$, hàm số nghịch biến nên $f(2) \leq f(x) \leq f(0)$, tức là $-8 \leq f(x) \leq 0$ (vì $f(2) = -2 \cdot 2^2 = -8$).

b) Kết hợp hai phần: với $-1 \leq x \leq 2$ luôn có $-8 \leq f(x) \leq 0$; giá trị $0$ đạt được tại $x = 0$ (đỉnh $O$ là điểm cao nhất của đồ thị) và giá trị $-8$ đạt được tại $x = 2$.

b) Vậy giá trị lớn nhất là $0$ khi $x = 0$, giá trị nhỏ nhất là $-8$ khi $x = 2$. Lưu ý: không thể chỉ tính giá trị tại hai đầu mút $x = -1$ và $x = 2$, vì hàm số không đồng biến hay nghịch biến trên cả đoạn $-1 \leq x \leq 2$.

Đáp số: a) $\boxed{\text{GTLN} = -2 \text{ khi } x = 1;\ \text{GTNN} = -18 \text{ khi } x = 3}$; b) $\boxed{\text{GTLN} = 0 \text{ khi } x = 0;\ \text{GTNN} = -8 \text{ khi } x = 2}$
\end{solution}
\end{problem}

% ---------------------------------------------------------------- đường thẳng qua một điểm, song song đường cho trước; giao điểm vô tỉ
\begin{problem}[9DS-HSB2-026]
% nguon: dummy (tự sinh để test)
% lop: 9
% kc_du_kien: 9DS-HSB2-03, 8DS-HSBN-05, 9DS-PT2-02
% do_kho_du_kien: 3
Cho parabol $(P)\colon y = -x^2$ và điểm $M(1; 1)$.

a) Viết phương trình đường thẳng $(d)$ đi qua $M$ và song song với đường thẳng $y = 2x + 5$.

b) Tìm toạ độ các giao điểm của $(d)$ và $(P)$.

\begin{solution}
a) $(d)$ song song với đường thẳng $y = 2x + 5$ nên có cùng hệ số góc $2$ và tung độ gốc khác $5$: $(d)\colon y = 2x + b$ với $b \neq 5$.

a) $(d)$ đi qua $M(1; 1)$ nên $1 = 2 \cdot 1 + b$, suy ra $b = -1$ (thoả mãn $b \neq 5$). Vậy $(d)\colon y = 2x - 1$.

b) Phương trình hoành độ giao điểm của $(d)$ và $(P)$: $-x^2 = 2x - 1 \Leftrightarrow x^2 + 2x - 1 = 0$.

b) Dùng công thức nghiệm thu gọn với $b' = 1$: $\Delta' = 1^2 - 1 \cdot (-1) = 2 > 0$, nên phương trình có hai nghiệm phân biệt $x_1 = -1 + \sqrt{2}$ và $x_2 = -1 - \sqrt{2}$.

b) Tính tung độ từ $y = 2x - 1$: với $x_1 = -1 + \sqrt{2}$ thì $y_1 = 2\left(-1 + \sqrt{2}\right) - 1 = -3 + 2\sqrt{2}$; với $x_2 = -1 - \sqrt{2}$ thì $y_2 = 2\left(-1 - \sqrt{2}\right) - 1 = -3 - 2\sqrt{2}$.

b) Kiểm tra trên $(P)$ bằng hằng đẳng thức bình phương: $x_1^2 = 1 - 2\sqrt{2} + 2 = 3 - 2\sqrt{2}$ nên $-x_1^2 = -3 + 2\sqrt{2} = y_1$; $x_2^2 = 1 + 2\sqrt{2} + 2 = 3 + 2\sqrt{2}$ nên $-x_2^2 = -3 - 2\sqrt{2} = y_2$ (đúng).

Đáp số: a) $\boxed{(d)\colon y = 2x - 1}$; b) $\boxed{\left(-1 + \sqrt{2}; -3 + 2\sqrt{2}\right) \text{ và } \left(-1 - \sqrt{2}; -3 - 2\sqrt{2}\right)}$
\end{solution}
\end{problem}

% ---------------------------------------------------------------- điểm trên parabol cách đều hai trục, hai cách
\begin{problem}[9DS-HSB2-027]
% nguon: dummy (tự sinh để test)
% lop: 9
% kc_du_kien: 9DS-HSB2-03, 9DS-HSB2-01, 9DS-PT2-01
% do_kho_du_kien: 3
Cho parabol $(P)\colon y = \dfrac{1}{3}x^2$. Tìm các điểm thuộc $(P)$, khác gốc toạ độ $O$, cách đều hai trục toạ độ.

\begin{solution}[cách 1: dùng điều kiện điểm thuộc parabol]
Điểm $M$ thuộc $(P)$ có hoành độ $x_0$ thì có tung độ $y_0 = \dfrac{1}{3}x_0^2$, nên $M\left(x_0; \dfrac{1}{3}x_0^2\right)$; điều kiện $M$ khác $O$ là $x_0 \neq 0$.

Khoảng cách từ $M$ đến trục hoành là $|y_0|$, đến trục tung là $|x_0|$. $M$ cách đều hai trục khi và chỉ khi $|y_0| = |x_0|$, tức là $\dfrac{1}{3}x_0^2 = |x_0|$ (vì $\dfrac{1}{3}x_0^2 \geq 0$ nên $|y_0| = y_0$).

Vì $x_0^2 = |x_0|^2$ nên phương trình trở thành $\dfrac{1}{3}|x_0|^2 - |x_0| = 0 \Leftrightarrow |x_0|\left(\dfrac{1}{3}|x_0| - 1\right) = 0$. Do $x_0 \neq 0$ nên $|x_0| \neq 0$, suy ra $\dfrac{1}{3}|x_0| = 1$, tức là $|x_0| = 3$, hay $x_0 = 3$ hoặc $x_0 = -3$.

Khi đó $y_0 = \dfrac{1}{3} \cdot 3^2 = 3$. Vậy có hai điểm cần tìm là $(3; 3)$ và $(-3; 3)$.

Đáp số: $\boxed{(3; 3) \text{ và } (-3; 3)}$
\end{solution}

\begin{solution}[cách 2: giao điểm của parabol với hai đường phân giác]
Điểm $M(x_0; y_0)$ cách đều hai trục toạ độ khi và chỉ khi $|y_0| = |x_0|$, tức là $y_0 = x_0$ hoặc $y_0 = -x_0$; nghĩa là $M$ thuộc đường thẳng $y = x$ hoặc đường thẳng $y = -x$. Vậy các điểm cần tìm là các giao điểm khác $O$ của $(P)$ với hai đường thẳng này.

Giao điểm của $(P)$ và đường thẳng $y = x$: phương trình hoành độ $\dfrac{1}{3}x^2 = x \Leftrightarrow x\left(\dfrac{1}{3}x - 1\right) = 0$ (phương trình bậc hai khuyết hệ số tự do, đặt nhân tử chung), nên $x = 0$ hoặc $x = 3$. Loại $x = 0$ (ứng với điểm $O$); với $x = 3$ thì $y = 3$, được điểm $(3; 3)$.

Giao điểm của $(P)$ và đường thẳng $y = -x$: $\dfrac{1}{3}x^2 = -x \Leftrightarrow x\left(\dfrac{1}{3}x + 1\right) = 0$, nên $x = 0$ hoặc $x = -3$. Loại $x = 0$; với $x = -3$ thì $y = -(-3) = 3$, được điểm $(-3; 3)$.

Đáp số: $\boxed{(3; 3) \text{ và } (-3; 3)}$
\end{solution}
\end{problem}

% ---------------------------------------------------------------- suy ra dấu của a từ một bất đẳng thức, so sánh, giải bất đẳng thức
\begin{problem}[9DS-HSB2-028]
% nguon: dummy (tự sinh để test)
% lop: 9
% kc_du_kien: 9DS-HSB2-02, 9DS-CAN-03, 9DS-BPT-01
% do_kho_du_kien: 4
Cho hàm số $y = f(x) = ax^2$ ($a \neq 0$). Biết rằng $f(-3) > f(2)$.

a) Xác định dấu của $a$. Hàm số đồng biến, nghịch biến khi nào?

b) Không tính $a$, hãy so sánh $f\left(1 + \sqrt{2}\right)$ và $f\left(-2\sqrt{2}\right)$.

c) Tìm tất cả các giá trị của $x$ sao cho $f(x) < f(3)$.

\begin{solution}
a) Ta có $f(-3) = a \cdot (-3)^2 = 9a$ và $f(2) = a \cdot 2^2 = 4a$. Từ $f(-3) > f(2)$ suy ra $9a > 4a \Leftrightarrow 5a > 0 \Leftrightarrow a > 0$ (chia hai vế cho $5 > 0$, giữ nguyên chiều bất đẳng thức).

a) Vì $a > 0$ nên theo tính chất của hàm số $y = ax^2$, hàm số nghịch biến khi $x < 0$ và đồng biến khi $x > 0$.

b) Vì $\left(-2\sqrt{2}\right)^2 = \left(2\sqrt{2}\right)^2$ nên $f\left(-2\sqrt{2}\right) = f\left(2\sqrt{2}\right)$ (đồ thị nhận trục $Oy$ làm trục đối xứng). Do đó chỉ cần so sánh $f\left(1 + \sqrt{2}\right)$ với $f\left(2\sqrt{2}\right)$, trong đó $1 + \sqrt{2}$ và $2\sqrt{2}$ đều là số dương.

b) So sánh hai số dương $1 + \sqrt{2}$ và $2\sqrt{2}$ qua bình phương của chúng: $\left(1 + \sqrt{2}\right)^2 = 3 + 2\sqrt{2}$ và $\left(2\sqrt{2}\right)^2 = 8$. Ta có $3 + 2\sqrt{2} < 8 \Leftrightarrow 2\sqrt{2} < 5$; đưa thừa số vào trong dấu căn, $2\sqrt{2} = \sqrt{8} < \sqrt{25} = 5$ (đúng). Vậy $\left(1 + \sqrt{2}\right)^2 < \left(2\sqrt{2}\right)^2$, suy ra $0 < 1 + \sqrt{2} < 2\sqrt{2}$.

b) Hàm số đồng biến khi $x > 0$ (câu a), nên từ $0 < 1 + \sqrt{2} < 2\sqrt{2}$ suy ra $f\left(1 + \sqrt{2}\right) < f\left(2\sqrt{2}\right) = f\left(-2\sqrt{2}\right)$.

c) $f(x) < f(3) \Leftrightarrow ax^2 < 9a \Leftrightarrow x^2 < 9$ (chia hai vế cho $a > 0$, giữ nguyên chiều bất đẳng thức).

c) Vì $x^2 = |x|^2$ và $9 = 3^2$, với hai số không âm $|x|$ và $3$ ta có $|x|^2 < 3^2 \Leftrightarrow |x| < 3 \Leftrightarrow -3 < x < 3$.

Đáp số: a) $\boxed{a > 0;\ \text{nghịch biến khi } x < 0, \text{ đồng biến khi } x > 0}$; b) $\boxed{f\left(1 + \sqrt{2}\right) < f\left(-2\sqrt{2}\right)}$; c) $\boxed{-3 < x < 3}$
\end{solution}
\end{problem}

% ---------------------------------------------------------------- hệ số a là tam thức theo tham số: dấu của một tích
\begin{problem}[9DS-HSB2-029]
% nguon: dummy (tự sinh để test)
% lop: 9
% kc_du_kien: 9DS-HSB2-02, 9DS-HSB2-01, 9DS-BPT-08
% do_kho_du_kien: 4
Cho hàm số $y = (m^2 - 4m + 3)x^2$, với $m$ là tham số, $m \neq 1$ và $m \neq 3$.

a) Tìm $m$ để hàm số đồng biến khi $x > 0$.

b) Tìm các giá trị nguyên của $m$ để đồ thị của hàm số nằm phía dưới trục hoành. Với giá trị $m$ đó, viết công thức của hàm số và tính giá trị của hàm số tại $x = -3$.

\begin{solution}
a) Phân tích hệ số thành nhân tử (tách hạng tử $-4m = -m - 3m$): $m^2 - 4m + 3 = m(m - 1) - 3(m - 1) = (m - 1)(m - 3)$. Với $m \neq 1$, $m \neq 3$ thì hệ số này khác $0$, nên hàm số có dạng $y = ax^2$ với $a = (m - 1)(m - 3)$.

a) Theo tính chất của hàm số $y = ax^2$, hàm số đồng biến khi $x > 0$ khi và chỉ khi $a > 0$, tức là $(m - 1)(m - 3) > 0$; tích của hai thừa số lớn hơn $0$ khi và chỉ khi hai thừa số cùng dấu.

a) Trường hợp hai thừa số cùng dương: $m - 1 > 0$ và $m - 3 > 0$, tức là $m > 1$ và $m > 3$; nghiệm chung là $m > 3$.

a) Trường hợp hai thừa số cùng âm: $m - 1 < 0$ và $m - 3 < 0$, tức là $m < 1$ và $m < 3$; nghiệm chung là $m < 1$.

a) Vậy hàm số đồng biến khi $x > 0$ khi và chỉ khi $m < 1$ hoặc $m > 3$ (các giá trị này đều thoả mãn $m \neq 1$, $m \neq 3$).

b) Theo tính chất của hàm số $y = ax^2$, đồ thị nằm phía dưới trục hoành khi và chỉ khi $a < 0$, tức là $(m - 1)(m - 3) < 0$ (dùng phân tích ở câu a); tích của hai thừa số nhỏ hơn $0$ khi và chỉ khi hai thừa số trái dấu.

b) Trường hợp $m - 1 > 0$ và $m - 3 < 0$: tức là $m > 1$ và $m < 3$; nghiệm chung là $1 < m < 3$.

b) Trường hợp $m - 1 < 0$ và $m - 3 > 0$: tức là $m < 1$ và $m > 3$; hai bất phương trình này không có nghiệm chung vì không có số nào vừa nhỏ hơn $1$ vừa lớn hơn $3$.

b) Vậy đồ thị nằm phía dưới trục hoành khi và chỉ khi $1 < m < 3$. Vì $m$ là số nguyên nên $m = 2$ (thoả mãn $m \neq 1$, $m \neq 3$).

b) Với $m = 2$, hệ số $a = 2^2 - 4 \cdot 2 + 3 = -1$, nên hàm số là $y = -x^2$. Thay $x = -3$: $y = -(-3)^2 = -9$.

Đáp số: a) $\boxed{m < 1 \text{ hoặc } m > 3}$; b) $\boxed{m = 2;\ y = -x^2;\ \text{tại } x = -3 \text{ thì } y = -9}$
\end{solution}
\end{problem}

% ---------------------------------------------------------------- bài toán thực tế: điểm tiếp đất của vận động viên trượt tuyết
\begin{problem}[9DS-HSB2-030]
% nguon: dummy (tự sinh để test)
% lop: 9
% kc_du_kien: 9DS-HSB2-03, 9DS-PT2-07, 9DS-HSB2-01
% do_kho_du_kien: 4
Một vận động viên trượt tuyết rời bệ nhảy tại điểm $O$ theo phương nằm ngang. Chọn hệ trục toạ độ $Oxy$ có gốc $O$ tại mép bệ nhảy, trục $Ox$ nằm ngang theo hướng bay, trục $Oy$ thẳng đứng hướng lên (đơn vị trên các trục là mét). Khi đó đường bay của vận động viên là một phần của parabol $(P)\colon y = -\dfrac{1}{20}x^2$ (ứng với $x \geq 0$), còn mặt dốc tiếp đất nằm trên đường thẳng $(d)\colon y = -\dfrac{1}{2}x - 10$.

a) Tìm toạ độ điểm tiếp đất $T$ (giao điểm của đường bay với mặt dốc).

b) Gọi $C$ là giao điểm của mặt dốc với trục $Oy$. Tính độ dài đoạn dốc $CT$.

c) Một vận động viên khác cũng rời bệ tại $O$ theo phương nằm ngang, có đường bay là một phần của parabol $y = ax^2$ và tiếp đất tại điểm của mặt dốc có hoành độ $30$. Tìm $a$.

\begin{solution}
a) Gọi hoành độ điểm tiếp đất là $x$ (mét, $x > 0$ vì vận động viên bay về phía trước). Hoành độ này là nghiệm của phương trình hoành độ giao điểm $-\dfrac{1}{20}x^2 = -\dfrac{1}{2}x - 10$; nhân hai vế với $-20$ rồi chuyển vế: $x^2 - 10x - 200 = 0$.

a) Dùng công thức nghiệm thu gọn với $b' = -5$: $\Delta' = (-5)^2 - 1 \cdot (-200) = 225$, $\sqrt{\Delta'} = 15$, nên $x_1 = 5 + 15 = 20$ và $x_2 = 5 - 15 = -10$.

a) Đối chiếu điều kiện $x > 0$: loại $x_2 = -10$ (nằm phía sau bệ nhảy, ngoài phần đường bay), nhận $x_1 = 20$. Tung độ: $y = -\dfrac{1}{2} \cdot 20 - 10 = -20$ (kiểm tra trên $(P)$: $-\dfrac{1}{20} \cdot 20^2 = -20$). Vậy $T(20; -20)$: vận động viên tiếp đất cách bệ $20$ m theo phương ngang và thấp hơn mép bệ $20$ m.

b) Thay $x = 0$ vào phương trình của $(d)$ được $y = -10$, nên $C(0; -10)$.

b) Gọi $H(20; -10)$. Khi đó $CH$ nằm ngang với $CH = 20 - 0 = 20$ (m), $HT$ thẳng đứng với $HT = -10 - (-20) = 10$ (m), nên tam giác $CHT$ vuông tại $H$. Theo định lí Pythagore: $CT^2 = CH^2 + HT^2 = 400 + 100 = 500$, suy ra $CT = \sqrt{500} = 10\sqrt{5}$ (m).

c) Điểm tiếp đất mới thuộc mặt dốc và có hoành độ $30$ nên tung độ là $y = -\dfrac{1}{2} \cdot 30 - 10 = -25$; vậy điểm tiếp đất là $T'(30; -25)$.

c) $T'$ thuộc đường bay $y = ax^2$ nên $-25 = a \cdot 30^2 = 900a$, suy ra $a = -\dfrac{25}{900} = -\dfrac{1}{36}$ (thoả mãn $a < 0$, đường bay đi xuống).

Đáp số: a) $\boxed{T(20; -20)}$; b) $\boxed{CT = 10\sqrt{5} \text{ m}}$; c) $\boxed{a = -\dfrac{1}{36}}$
\end{solution}
\end{problem}

% ---------------------------------------------------------------- giao điểm, độ dài, diện tích; điểm trên cung AB cho diện tích lớn nhất
\begin{problem}[9DS-HSB2-031]
% nguon: dummy (tự sinh để test)
% lop: 9
% kc_du_kien: 9DS-HSB2-03, 8DS-HDT-01, 9DS-PT2-02
% do_kho_du_kien: 5
Cho parabol $(P)\colon y = x^2$ và đường thẳng $(d)\colon y = x + 6$.

a) Tìm toạ độ các giao điểm $A$, $B$ của $(d)$ và $(P)$ (với hoành độ của $A$ nhỏ hơn hoành độ của $B$).

b) Tính độ dài đoạn $AB$ và diện tích tam giác $OAB$ ($O$ là gốc toạ độ).

c) Điểm $C$ thuộc $(P)$ có hoành độ $c$ với $-2 < c < 3$. Tìm toạ độ của $C$ để tam giác $ABC$ có diện tích lớn nhất và tính diện tích lớn nhất đó.

\begin{solution}
a) Phương trình hoành độ giao điểm của $(P)$ và $(d)$: $x^2 = x + 6 \Leftrightarrow x^2 - x - 6 = 0$; có $\Delta = (-1)^2 - 4 \cdot 1 \cdot (-6) = 25$, $\sqrt{\Delta} = 5$, nên $x = \dfrac{1 + 5}{2} = 3$ hoặc $x = \dfrac{1 - 5}{2} = -2$.

a) Tính tung độ từ $y = x^2$: với $x = -2$ thì $y = 4$; với $x = 3$ thì $y = 9$ (thử lại bằng $y = x + 6$: $-2 + 6 = 4$, $3 + 6 = 9$). Vậy $A(-2; 4)$ và $B(3; 9)$.

b) Gọi $K(3; 4)$; khi đó $AK$ nằm ngang với $AK = 3 - (-2) = 5$, $KB$ thẳng đứng với $KB = 9 - 4 = 5$, tam giác $AKB$ vuông tại $K$. Theo định lí Pythagore: $AB^2 = 5^2 + 5^2 = 50$, suy ra $AB = 5\sqrt{2}$.

b) $(d)$ cắt trục tung tại $E(0; 6)$, $OE = 6$. Vì $A$, $B$ có hoành độ trái dấu nên nằm về hai phía của trục tung, $E$ nằm trên đoạn $AB$ và chia tam giác $OAB$ thành hai tam giác $OAE$, $OBE$ có chung cạnh $OE$ nằm trên trục tung; đường cao hạ từ $A$, $B$ xuống $OE$ lần lượt bằng $|-2| = 2$ và $|3| = 3$.

b) Do đó $S_{OAB} = S_{OAE} + S_{OBE} = \dfrac{1}{2} \cdot 6 \cdot 2 + \dfrac{1}{2} \cdot 6 \cdot 3 = 6 + 9 = 15$ (đơn vị diện tích).

c) Điểm $C$ thuộc $(P)$ nên $C(c; c^2)$. Đường thẳng đứng $x = c$ cắt $(d)$ tại $D(c; c + 6)$; vì $-2 < c < 3$ nên $D$ nằm trên đoạn $AB$ (giữa $A$ và $B$).

c) Độ dài $CD = (c + 6) - c^2 = -(c^2 - c - 6) = -(c + 2)(c - 3)$; với $-2 < c < 3$ thì $c + 2 > 0$ và $c - 3 < 0$, nên $CD > 0$, tức là $D$ nằm phía trên $C$ và $CD = 6 + c - c^2$.

c) Đoạn $CD$ chia tam giác $ABC$ thành hai tam giác $ACD$ và $BCD$ có chung cạnh $CD$ thẳng đứng; đường cao hạ từ $A$ và $B$ xuống đường thẳng $CD$ lần lượt bằng $c - (-2) = c + 2$ và $3 - c$. Do đó $S_{ABC} = \dfrac{1}{2} \cdot CD \cdot (c + 2) + \dfrac{1}{2} \cdot CD \cdot (3 - c) = \dfrac{5}{2} \cdot CD = \dfrac{5}{2}\left(6 + c - c^2\right)$.

c) Biến đổi bằng hằng đẳng thức bình phương của một hiệu: $6 + c - c^2 = \dfrac{25}{4} - \left(c^2 - c + \dfrac{1}{4}\right) = \dfrac{25}{4} - \left(c - \dfrac{1}{2}\right)^2 \leq \dfrac{25}{4}$, dấu bằng xảy ra khi và chỉ khi $c = \dfrac{1}{2}$ (thoả mãn $-2 < c < 3$).

c) Vậy $S_{ABC} \leq \dfrac{5}{2} \cdot \dfrac{25}{4} = \dfrac{125}{8}$; diện tích lớn nhất bằng $\dfrac{125}{8}$ khi $c = \dfrac{1}{2}$, lúc đó $C\left(\dfrac{1}{2}; \dfrac{1}{4}\right)$. (Kiểm tra với $c = 0$, tức $C$ trùng $O$: $S = \dfrac{5}{2} \cdot 6 = 15$, khớp với câu b.)

Đáp số: a) $\boxed{A(-2; 4),\ B(3; 9)}$; b) $\boxed{AB = 5\sqrt{2},\ S_{OAB} = 15 \text{ (đơn vị diện tích)}}$; c) $\boxed{C\left(\dfrac{1}{2}; \dfrac{1}{4}\right),\ S_{\max} = \dfrac{125}{8} \text{ (đơn vị diện tích)}}$
\end{solution}
\end{problem}

\end{document}
